\cleardoublepage % memastikan bab baru dimulai di halaman ganjil (kanan)
\chapter{ORGANISASI ATAU LINGKUNGAN KERJA PRAKTIK}
\label{chap:organisasi_lingkungan}
\section{Profil Perusahaan}
\label{sec:profil_perusahaan}
PT Paragon Technology and Innovation merupakan perusahaan kecantikan asal Indonesia yang beroperasi di bawah Paragon Corp.
Perusahaan ini bergerak di bidang manufaktur, pengembangan, dan distribusi produk kecantikan untuk pasar domestik maupun internasional.
Produk perusahaan mencakup berbagai merek kosmetik dan perawatan tubuh yang telah dikenal luas di Indonesia.

Tujuan (\textit{purpose}) perusahaan adalah menciptakan kebaikan yang lebih besar bagi masyarakat melalui inovasi.
Orientasi tersebut tercermin dalam komitmen perusahaan untuk menghadirkan produk berkualitas tinggi sekaligus memberi manfaat bagi karyawan (\textit{Paragonians}), mitra bisnis, masyarakat, dan lingkungan.

Visi perusahaan dirumuskan sebagai berikut.
\begin{quote}
    Menjadi perusahaan yang berkomitmen pada tata kelola perusahaan terbaik yang dilengkapi dengan perbaikan berkelanjutan, agar setiap hari lebih baik dari hari sebelumnya, melalui produk berkualitas tinggi yang bermanfaat bagi \textit{Paragonians}, mitra, masyarakat, dan lingkungan.
\end{quote}

Untuk mewujudkan visi tersebut, perusahaan menetapkan tujuh misi, yaitu sebagai berikut.
\begin{enumerate}
    \item \textbf{Developing Paragonians}: terus belajar dan mengembangkan \textit{Paragonians} yang kompeten dengan keunggulan bersaing;
    \item \textbf{Creating Kindness for Consumers}: mendengarkan kebutuhan konsumen dan menciptakan produk yang melampaui harapan konsumen;
    \item \textbf{Continuous Improvement}: terus berinovasi dalam setiap proses guna meningkatkan kualitas produk;
    \item \textbf{Grow Together}: bekerja sama dengan mitra bisnis demi keuntungan bersama;
    \item \textbf{Preserving the Earth}: berupaya melindungi bumi secara berkelanjutan sebagai sumber kehidupan;
    \item \textbf{Supporting Education and Health of the Nation}: mendukung pengembangan generasi baru melalui peningkatan sumber daya manusia agar berpengetahuan, berpengalaman, dan sehat;
    \item \textbf{Developing Business}: memperluas ragam produk dan layanan melalui ekspansi bisnis.
\end{enumerate}

Dalam industri kecantikan nasional, Paragon termasuk perusahaan terbesar dengan skala operasional yang luas.
Posisi tersebut didukung oleh portofolio merek yang kuat, jaringan distribusi yang menjangkau berbagai segmen pasar, serta komitmen terhadap inovasi produk dan proses bisnis.
Dari perspektif kerja praktik penulis, lingkungan perusahaan yang mendorong perbaikan berkelanjutan dan inovasi relevan dengan pengembangan sistem informasi berbasis akal imitasi (AI) untuk mendukung efisiensi operasional perusahaan.
\section{Struktur Organisasi}
\label{sec:struktur_organisasi}
\textbf{JANGAN LUPA MINTA MAS TEGAR}
PT Paragon Technology and Innovation memiliki struktur organisasi berjenjang yang mengelompokkan fungsi bisnis dan teknologi informasi ke dalam direktorat, departemen, serta divisi.
Selama kerja praktik, penulis berada pada Direktorat Enterprise Technology and Data Excellence, yang bertanggung jawab atas pengembangan dan penguatan kapabilitas teknologi informasi serta data di lingkungan perusahaan.
Di bawah direktorat tersebut terdapat Departemen Software Engineering yang menangani perancangan, pengembangan, dan pemeliharaan perangkat lunak pendukung proses bisnis perusahaan.

Penulis ditempatkan pada Divisi \label{abbr:FATL}FATL (\textit{Finance, Accounting, Tax and Legal}), yaitu divisi yang menangani sistem informasi terkait keuangan, akuntansi, perpajakan, dan aspek legal perusahaan.
Posisi penulis dalam struktur organisasi dapat digambarkan sebagai berikut: Direktorat Enterprise Technology and Data Excellence $\rightarrow$ Departemen Software Engineering $\rightarrow$ Divisi FATL $\rightarrow$ bagian \textit{Web Tukar Faktur} $\rightarrow$ penulis (\textit{software engineer intern}).
\section{Unit Kerja}
\label{sec:unit_kerja}
Unit kerja tempat penulis melaksanakan kerja praktik adalah bagian \textit{Web Tukar Faktur} pada Divisi FATL.
Bagian ini bertanggung jawab mengembangkan dan memelihara \textit{Web Tukar Faktur}, sistem informasi yang mendukung proses pertukaran faktur antara perusahaan dan pemasok.
Melalui sistem ini, proses verifikasi, koreksi, dan penyelesaian administrasi faktur dapat dilakukan secara terintegrasi sehingga mendukung kepatuhan perpajakan dan kelancaran hubungan dengan mitra bisnis.
Selain pengembangan fitur baru, unit kerja ini juga bertanggung jawab melakukan \textit{bug fixing} pada operasional sehari-hari guna menjaga kestabilan sistem.
Apabila muncul keluhan atau kendala dari pengguna, tim bertugas menerima laporan, menganalisis penyebab masalah, serta menindaklanjuti perbaikan agar layanan \textit{Web Tukar Faktur} tetap dapat digunakan dengan andal.

Tim pengembang pada bagian \textit{Web Tukar Faktur} terdiri atas penulis sebagai peserta magang dan satu karyawan tetap yang bertindak sebagai mentor penulis selama pelaksanaan kerja praktik.
Selama periode kerja praktik, unit kerja ini menjalankan inisiatif pengembangan subsistem Sistem Koreksi Tarif Pajak Penghasilan Berbasis Akal Imitasi (AI) guna mengatasi ketidaksesuaian tarif PPh pada proses tukar faktur.
Kontribusi unit kerja ini sejalan dengan fungsi Divisi FATL dalam mendukung operasional keuangan dan kepatuhan perpajakan perusahaan.
\section{Lingkungan Kerja} 
\label{sec:lingkungan_kerja}
Lingkungan kerja di PT Paragon Technology and Innovation ditandai oleh budaya yang kental dengan sifat kekeluargaan.
Tempat kerja penulis berada di Gedung Nomor 3, yaitu gedung bekas rumah pendiri Paragon yang kini dialihfungsikan sebagai ruang kerja.
Lingkungan fisik tersebut memberi kesan hangat dan akrab, sehingga aktivitas kerja sehari-hari terasa lebih personal dan tidak sekaku lingkungan perkantoran konvensional.

Suasana kerja yang demikian mendorong kenyamanan pegawai dalam bekerja sekaligus berkomunikasi satu sama lain.
Interaksi antarkaryawan berlangsung terbuka tanpa memandang tingkat hierarki perusahaan, sehingga penulis sebagai peserta magang pun dapat berdiskusi, bertanya, maupun menyampaikan ide secara langsung kepada rekan kerja maupun mentor.
Budaya ini selaras dengan prinsip Paragon yang lebih menghargai kontribusi ketimbang sekadar posisi atau jenjang jabatan.
Akibatnya, setiap karyawan---termasuk penulis---didorong untuk aktif berkontribusi.

Dari sisi pelaksanaan teknis, pegawai diberikan kebebasan mikro dalam pekerjaan sehari-hari, asalkan tetap mematuhi acuan resmi perusahaan.
Contoh kebebasan tersebut meliputi pemilihan \textit{library}, metode implementasi, serta \textit{design pattern} yang digunakan dalam pengembangan perangkat lunak.
Fleksibilitas ini memberi ruang eksplorasi yang diperlukan penulis, khususnya pada tahap perancangan \textit{proof of concept} (POC) dan pengembangan \textit{backend} sistem AI, tanpa mengabaikan standar teknis yang berlaku di unit kerja.

Selama pelaksanaan kerja praktik, penulis beradaptasi dengan lingkungan kerja tersebut melalui partisipasi aktif dalam diskusi tim, pengembangan perangkat lunak, serta kolaborasi langsung dengan mentor.
Lingkungan yang inklusif dan meritokratis tersebut membantu penulis memperoleh pengalaman profesional yang mendekati praktik kerja nyata di industri, sekaligus memperkuat kemampuan komunikasi dan pemecahan masalah dalam konteks proyek \textit{Web Tukar Faktur}.








